\documentclass[border=4pt]{standalone}
\usepackage[siunitx]{circuitikz}
\usetikzlibrary{calc}

\begin{document}
\begin{circuitikz}[font=\small,line width=0.9pt]
  \ctikzset{tripoles/thickness=1,logic ports/thickness=1,
    flipflops/thickness=1}
  \node[font=\bfseries] at (6.7,9.0)
    {(a) Four-bit parallel-in/parallel-out register};
  \foreach \i/\xx in {0/1.3,1/4.7,2/8.1,3/11.5} {
    \node[flipflop D,scale=0.82] (P\i) at (\xx,6.8) {};
    \node[above,font=\scriptsize] at (P\i.north) {stage $\i$};
    \draw (\xx-1.15,8.2) node[above] {$D_{\i}$} |- (P\i.pin 1);
    \draw (P\i.pin 6) -- ++(0.5,0) node[right] {$Q_{\i}$};
  }
  \draw (-0.15,4.75) node[left] {$CLK$} -- (12.4,4.75);
  \foreach \i/\xx in {0/1.3,1/4.7,2/8.1,3/11.5} {
    \coordinate (PC\i) at ($(P\i.pin 3)+(-0.35,0)$);
    \coordinate (PB\i) at (PC\i |- 0,4.75);
    \draw (PB\i) -- (PC\i) -- (P\i.pin 3);
    \fill (PB\i) circle (1.4pt);
  }
  \node[align=center] at (6.7,4.05)
    {Independent D inputs are sampled together by one rising-edge clock.};

  \draw[densely dashed] (-0.2,3.45) -- (13.5,3.45);
  \node[font=\bfseries] at (6.7,2.9)
    {(b) Four-stage serial shift register with parallel Q taps};
  \foreach \i/\xx in {0/1.3,1/4.7,2/8.1,3/11.5} {
    \node[flipflop D,scale=0.82] (S\i) at (\xx,0.75) {};
    \node[above,font=\scriptsize] at (S\i.north) {stage $\i$};
  }
  \draw ($(S0.pin 1)+(-0.9,0)$) node[left] {$SER_{in}$} -- (S0.pin 1);
  \foreach \i/\j in {0/1,1/2,2/3} {
    \coordinate (J\i) at ($(S\i.pin 6)+(0.25,0)$);
    \draw (S\i.pin 6) -- (J\i) -- (S\j.pin 1);
    \fill (J\i) circle (1.4pt);
    \draw (J\i) -- ++(0,0.75) node[above] {$Q_{\i}$ tap};
  }
  \coordinate (J3) at ($(S3.pin 6)+(0.25,0)$);
  \draw (S3.pin 6) -- (J3);
  \fill (J3) circle (1.4pt);
  \draw (J3) -- ++(0,0.75) node[above] {$Q_3$ tap};
  \draw (J3) -- ++(1.0,0) node[right] {$SER_{out}$};

  \draw (-0.15,-1.35) node[left] {$CLK$} -- (12.4,-1.35);
  \foreach \i/\xx in {0/1.3,1/4.7,2/8.1,3/11.5} {
    \coordinate (SC\i) at ($(S\i.pin 3)+(-0.35,0)$);
    \coordinate (SB\i) at (SC\i |- 0,-1.35);
    \draw (SB\i) -- (SC\i) -- (S\i.pin 3);
    \fill (SB\i) circle (1.4pt);
  }
  \node[align=center] at (6.7,-2.05)
    {Each edge moves one old Q value into the D input of the next stage.};
\end{circuitikz}
\end{document}
